%! Properties and Applications of Determinants — Unit 5, Lesson 5.2 — Slides (Beamer)
\documentclass[aspectratio=169, 11pt]{beamer}
\usepackage{linalg-beamer}   % provides \burgundyheader{} and \sectionlabel[color]{}
\newcommand{\T}{\mathsf{T}}
\newcommand{\xx}{\mathbf{x}}

\begin{document}

% TITLE SLIDE (CourseName is not defined in beamer, so the name is literal)
{
\setbeamercolor{background canvas}{bg=burgundy}
\begin{frame}[plain]
  \vspace{2.8cm}\hspace{0.55cm}%
  \begin{minipage}{0.9\textwidth}
    {\color{goldacc}\bfseries\small LESSON 5.2}\\[0.5cm]
    {\color{white}\bfseries\LARGE Properties and Applications of Determinants}\\[1.0cm]
    {\color{blushmid}\small Linear Algebra~$\cdot$~Unit 5: Determinants and Linear Transformations}
  \end{minipage}
\end{frame}
}

% HOOK
\begin{frame}[t, plain]
  \burgundyheader{Is there a faster way than cofactors?}
  \sectionlabel{Hook}
  \begin{columns}[T]
    \begin{column}{0.56\textwidth}
      Expanding a $3\times3$ by cofactors is \emph{three} $2\times2$s and a lot of bookkeeping.\\[8pt]
      But you already own a machine that makes any matrix \textbf{triangular} --- \textbf{elimination}
      (Unit~2) --- and a triangular determinant is free.\\[8pt]
      The missing link: \emph{does elimination change the determinant?}
    \end{column}
    \begin{column}{0.42\textwidth}
      \centering
      \begin{tikzpicture}[scale=0.9]
        \fill[burgundy!12] (0,0)--(1.6,0)--(1.6,1.3)--(0,1.3)--cycle;
        \draw[charcoal, thin] (0,0)--(1.6,0)--(1.6,1.3)--(0,1.3)--cycle;
        \draw[->, very thick, burgundy!70] (2.0,0.65)--(2.9,0.65) node[midway,above]{\scriptsize shear};
        \fill[greenacc!16] (3.3,0)--(4.9,0)--(5.6,1.3)--(4.0,1.3)--cycle;
        \draw[charcoal, thin] (3.3,0)--(4.9,0)--(5.6,1.3)--(4.0,1.3)--cycle;
        \node[charcoal] at (0.8,0.6){\scriptsize $bh$};
        \node[charcoal] at (4.5,0.6){\scriptsize $bh$};
      \end{tikzpicture}\\[4pt]
      {\scriptsize A shear keeps base \& height.}
    \end{column}
  \end{columns}
  \vspace{2pt}
  \begin{block}{Warm-Up punchline}
    On $M=\begin{bsmallmatrix} 1 & 2\\ 2 & 5 \end{bsmallmatrix}$, one elimination step gave $\det M=\det U=1$.
    \textbf{Elimination did not change the determinant.}
  \end{block}
\end{frame}

% THE RULES + SHEAR
\begin{frame}[t, plain]
  \burgundyheader{The rules --- each one is about volume}
  \sectionlabel[goldacc]{Rules}
  \begin{itemize}
    \setlength{\itemsep}{5pt}
    \item $\det I = 1$ \; (the unit cube has volume $1$).
    \item \textbf{Swap} two rows/columns $\Rightarrow$ $\det$ is \textbf{negated} (a mirror reflection).
    \item \textbf{Scale} a row by $t$ $\Rightarrow$ $\det$ is \textbf{multiplied} by $t$ (stretch an edge).
    \item \textbf{Shear} --- add a multiple of one row to another $\Rightarrow$ $\det$ is \textbf{unchanged}
          (slide the box; base \& height stay).
  \end{itemize}
  \vspace{2pt}
  \begin{block}{Check}
    $\det\begin{bsmallmatrix} 2 & 1\\ 1 & 3 \end{bsmallmatrix}=5$, and after a shear
    $\det\begin{bsmallmatrix} 2 & 0\\ 1 & 2.5 \end{bsmallmatrix}=5.$ \; Two equal rows $\Rightarrow\det=0$ (flat).
  \end{block}
\end{frame}

% PRODUCT OF PIVOTS
\begin{frame}[t, plain]
  \burgundyheader{Fast determinant: product of the pivots}
  \sectionlabel{Method}
  \begin{columns}[T]
    \begin{column}{0.5\textwidth}
      Shears don't change $\det$; swaps only flip its sign. So reduce to triangular and
      \textbf{multiply the pivots}:
      \[ \det A = \pm(\text{product of pivots}). \]
      One sign flip per row swap.
    \end{column}
    \begin{column}{0.5\textwidth}
      \centering
      $\begin{bsmallmatrix} 1 & 2 & 0\\ 2 & 5 & 1\\ 0 & 1 & 3 \end{bsmallmatrix}
       \xrightarrow{\;R_2-2R_1\;}\xrightarrow{\;R_3-R_2\;}
       \begin{bsmallmatrix} 1 & 2 & 0\\ 0 & 1 & 1\\ 0 & 0 & 2 \end{bsmallmatrix}$\\[8pt]
      pivots $1,1,2 \Rightarrow \det = 2.$\\[4pt]
      {\scriptsize No cofactors, no swaps.}
    \end{column}
  \end{columns}
\end{frame}

% PRODUCT RULE + ROAD AHEAD
\begin{frame}[t, plain]
  \burgundyheader{Determinants multiply --- and what it buys}
  \sectionlabel[goldacc]{Product rule}
  \begin{itemize}
    \setlength{\itemsep}{6pt}
    \item \textbf{Product rule:} $\det(AB)=\det A\,\det B$ (compose maps $\Rightarrow$ multiply volume factors).
          E.g.\ $\det A=5,\ \det B=6\Rightarrow\det(AB)=30$.
    \item \textbf{Inverses:} $\det A\cdot\det A^{-1}=\det I=1$, so $\det A^{-1}=1/\det A$.
    \item \textbf{Singular:} $\det A=0\ \Leftrightarrow\ A$ has \textbf{no inverse}. \; And $\det A^{\T}=\det A$.
  \end{itemize}
  \vspace{2pt}
  \begin{block}{The road ahead}
    \textbf{5.3} linear transformations: $|\det A|$ is the factor every area/volume is multiplied by,
    and its \emph{sign} tells you if orientation was kept or flipped.
  \end{block}
\end{frame}

\end{document}
